This report evaluated exact-quota mixed-precision FIR quantization as
an instance of unit-capacity transportation. The DP implements a known
cost-constrained optimization principle; no algorithmic novelty is
claimed. Its coefficient-SSE optimum was verified by 29 exhaustive
checks and 108 MILP comparisons (maximum discrepancy
$4.34\times10^{-19}$).

Hardware evidence comprises six routed 8-tap runs and 942 bit-exact
board comparisons, plus four routed 48-tap runs and 14,888 RTL output
comparisons. In the 48-tap base case, DP meets the stopband limit where
same-quota greedy fails, at the cost of more LUTs and registers.
Across five designs and 15 design/quota combinations, DP meets both
response limits in two cases and greedy in none; greedy nevertheless
has a lower stopband peak in some failed cases. Coefficient-SSE
optimality therefore does not guarantee spectral compliance.

The band-weighted response objective contains cross-terms, while its
uniform diagonal approximation leaves the SSE allocation unchanged.
The 48-tap binary-product MILP had a weak root relaxation and was
manually stopped after about 93 seconds under the pre-declared
slow-progress condition. No solver exit flag, final allocation, or
optimality certificate was returned. This observed run does not
establish mathematical intractability.

Future work includes tighter response-aware formulations or dedicated
MIQP solvers, and 48-tap board testing, power measurement, and Fmax
characterization. The released code and data preserve both the positive
validation results and these limitations.
